% from ega3-1-fr.tex
\oldpage[0\textsubscript{III}]{73}
s'identifie canoniquement à
$\mathrm E_\infty^{pq}(\mathbf X^\bullet)=
\mathrm Z_\infty^{pq}(\mathbf X^\bullet)/
\mathrm B_\infty^{pq}(\mathbf X^\bullet)$. On remarquera que les objets
$\mathrm Z_r^{pq}(\mathbf X^\bullet)$,
$\mathrm B_r^{pq}(\mathbf X^\bullet)$,
$\mathrm E_r^{pq}(\mathbf X^\bullet)$ $(1\leq r\leq+\infty)$ et
$d_r^{pq}$ dépendent \emph{fonctoriellement} du système projectif
$\mathbf X^\bullet$ soumis aux restrictions de
(\hyperref[0.13.5.5-fr]{13.5.5}), et que les morphismes définis dans
(\hyperref[0.13.5.5-fr]{13.5.5}) et
(\hyperref[0.13.5.6-fr]{13.5.6}) sont fonctoriels.
\end{env}

\subsection{Suite spectrale d'un foncteur relative à un objet muni d'une filtration finie}
\label{subsection:0.13.6-fr}

\begin{env}[13.6.1]
\label{0.13.6.1-fr}
Soient $C$, $C'$ deux catégories abéliennes, $\mathrm T:C\to C'$ un
foncteur additif covariant. Supposons que tout objet de $C$ soit isomorphe
à un sous-objet d'un objet injectif, de sorte que les foncteurs dérivés
droits $\mathrm R^p\mathrm T$ $(p\geq0)$ existent.
\end{env}

\begin{lemma}[13.6.2]
\label{0.13.6.2-fr}
\emph{Soit $A$ un objet de $C$, muni d'une filtration finie
$(\mathrm F^i(A))_{i\in\mathbb Z}$. Il existe une résolution injective
$X^\bullet=(X^j)_{j\geq0}$ de $A$ munie d'une filtration finie
$(\mathrm F^i(X^\bullet))_{i\in\mathbb Z}$ telle que la relation
$\mathrm F^i(A)=A$ (resp. $\mathrm F^i(A)=0$) entraîne
$\mathrm F^i(X^\bullet)=X^\bullet$ (resp. $\mathrm F^i(X^\bullet)=0$) et
que, pour tout $i\in\mathbb Z$, $\mathrm F^i(X^\bullet)$ soit une
résolution injective de $\mathrm F^i(A)$.}
\end{lemma}
\begin{proof}
Soit $p$ (resp. $q>p$) le plus grand indice tel que
$\mathrm F^p(A)=A$ (resp. le plus petit indice pour lequel
$\mathrm F^q(A)=0$). On raisonne par récurrence sur $q-p$, le lemme étant
évident pour $q-p=1$. Ayant formé une résolution injective
${X''}^\bullet$ de $A/\mathrm F^{q-1}(A)$ ayant les propriétés voulues,
on considère la suite exacte
\[
0\to\mathrm F^{q-1}(A)\to A\to A/\mathrm F^{q-1}(A)\to0,
\]
on prend une résolution injective ${X'}^\bullet$ de
$\mathrm F^{q-1}(A)$, puis on détermine une résolution injective
$X^\bullet$ de $A$ de façon à avoir une suite exacte
\[
0\to {X'}^\bullet\to X^\bullet\to {X''}^\bullet\to0
\]
compatible avec la précédente (M, V, 2.2)~; il est clair que
$X^\bullet$ répond à la question.
\end{proof}

\begin{corollary}[13.6.3]
\label{0.13.6.3-fr}
\emph{Soit $B$ un second objet de $C$, muni d'une filtration finie
$(\mathrm F^i(B))_{i\in\mathbb Z}$, $s$ un entier, et soit $u:A\to B$
un morphisme tel que $u(\mathrm F^i(A))\subset\mathrm F^{i+s}(B)$ pour
tout $i\in\mathbb Z$. Si $Y^\bullet=(Y^j)_{j\geq0}$ est une résolution
injective de $B$ munie d'une filtration
$(\mathrm F^i(Y^\bullet))_{i\in\mathbb Z}$ ayant les propriétés énoncées
en (\hyperref[0.13.6.2-fr]{13.6.2}), il existe un morphisme
$v:X^\bullet\to Y^\bullet$ compatible avec $u$ et tel que
$v(\mathrm F^i(X^\bullet))\subset\mathrm F^{i+s}(Y^\bullet)$ pour tout
$i\in\mathbb Z$. En outre, deux tels morphismes $v$, $v'$ sont
homotopes.}
\end{corollary}

Cela résulte aussitôt par récurrence sur $q-p$ de la construction
précédente et de (M, V, 2.3).

\begin{env}[13.6.4]
\label{0.13.6.4-fr}
Sous les hypothèses de (\hyperref[0.13.6.2-fr]{13.6.2}), considérons
maintenant le complexe $\mathrm T(X^\bullet)$ dans $C'$, qui est
évidemment filtré par les complexes
$\mathrm T(\mathrm F^i(X^\bullet))$, puisque
$\mathrm F^i(X^\bullet)$ est facteur direct de $X^\bullet$. Il résulte de
(\hyperref[0.13.6.3-fr]{13.6.3}) que la \emph{suite spectrale} de ce
complexe filtré ne dépend que de l'objet filtré $A$, à un isomorphisme
près. Son aboutissement est la cohomologie $\mathrm R^\bullet\mathrm T(A)$,
avec la filtration
\begin{equation}
\label{0.13.6.4.1-fr}
\tag{13.6.4.1}
\begin{split}
\mathrm F^p(\mathrm R^n\mathrm T(A))
&=\operatorname{Im}\bigl(\mathrm R^n\mathrm T(\mathrm F^p(A))
\longrightarrow\mathrm R^n\mathrm T(A)\bigr)\\
&=\operatorname{Ker}\bigl(\mathrm R^n\mathrm T(A)
\longrightarrow\mathrm R^n\mathrm T(A/\mathrm F^p(A))\bigr)
\end{split}
\end{equation}
(\hyperref[0.11.2.2-fr]{11.2.2}), et son terme $\mathrm E_1$ est donné
par
\begin{equation}
\label{0.13.6.4.2-fr}
\tag{13.6.4.2}
\mathrm E_1^{pq}=\mathrm R^{p+q}\mathrm T(\operatorname{gr}^p(A))
\end{equation}
$\operatorname{gr}^p(A)$ désignant comme d'ordinaire
$\mathrm F^p(A)/\mathrm F^{p+1}(A)$. Il est clair, d'après
(\hyperref[0.11.2.2-fr]{11.2.2}), que la filtration de l'aboutissement
est finie, et que pour $p$, $q$ donnés, les suites des
